% ============================================================
%  IMR GUI — Quick-Start Guide
%  Compile with: pdflatex IMR_GUI_quickstart.tex  (twice)
%  Packages required: geometry, graphicx, xcolor, tcolorbox,
%                     booktabs, hyperref, amsmath, enumitem
% ============================================================
\documentclass[11pt, letterpaper]{article}

\usepackage[margin=1in]{geometry}
\usepackage{graphicx}
\usepackage{xcolor}
\usepackage[most]{tcolorbox}
\usepackage{booktabs}
\usepackage{amsmath}
\usepackage{enumitem}
\usepackage{hyperref}
\usepackage{float}
\usepackage{microtype}
\usepackage{parskip}

% --- Color palette ---
\definecolor{noteblue}{RGB}{220, 235, 255}
\definecolor{warnamber}{RGB}{255, 243, 205}
\definecolor{okgreen}{RGB}{212, 237, 218}
\definecolor{accentgray}{RGB}{90, 90, 90}

% --- Callout boxes ---
\tcbset{
  notebox/.style={
    colback=noteblue, colframe=blue!40!black,
    fonttitle=\bfseries\small, title={Note}, left=4pt, right=4pt,
    top=3pt, bottom=3pt, boxrule=0.5pt
  },
  warnbox/.style={
    colback=warnamber, colframe=orange!70!black,
    fonttitle=\bfseries\small, title={Important}, left=4pt, right=4pt,
    top=3pt, bottom=3pt, boxrule=0.5pt
  },
  okbox/.style={
    colback=okgreen, colframe=green!40!black,
    fonttitle=\bfseries\small, title={Success Criterion}, left=4pt, right=4pt,
    top=3pt, bottom=3pt, boxrule=0.5pt
  },
}

% --- Step counter (integer steps) ---
\newcounter{stepctr}
\newcommand{\step}[1]{%
  \stepcounter{stepctr}%
  \vspace{6pt}
  \noindent
  {\large\bfseries\color{blue!60!black} Step~\thestepctr:\quad #1}
  \vspace{2pt}\par
}

% --- Half-step (e.g. 4.5): prints current counter value + ".5", no increment ---
\newcommand{\halfstep}[1]{%
  \vspace{6pt}
  \noindent
  {\large\bfseries\color{blue!60!black} Step~\thestepctr.5:\quad #1}
  \vspace{2pt}\par
}

\hypersetup{colorlinks=true, linkcolor=blue!60!black, urlcolor=blue!60!black}

% ============================================================
\begin{document}

\begin{center}
  {\LARGE\bfseries Quick-Start Guide}\\[4pt]
  {\large IMR Bubble Dynamics Fitting --- IMR GUI}\\[2pt]
  {\small\color{accentgray} For first-time users. No programming required.}
\end{center}

\vspace{4pt}
\hrule
\vspace{8pt}

% ---- Overview -----------------------------------------------
\section*{Overview}

This guide walks you through fitting a \textbf{single-bubble inertial
microrheology (IMR)} dataset to extract the shear modulus $G$ and
viscosity $\mu$ of a soft material using the \textbf{NHKV\,(Rmax) model}.
Follow every step in order. If anything looks unexpected, ask your TA
before changing settings not covered here.

% ---- Step 1 -------------------------------------------------
\step{Prepare Your Experimental Data File}

Your data must be saved as a \texttt{.mat} file containing two
one-dimensional arrays:

\begin{center}
\begin{tabular}{lll}
\toprule
\textbf{Quantity} & \textbf{Variable name} & \textbf{Unit} \\
\midrule
Bubble radius & \texttt{R} & meters (m) \\
Time          & \texttt{t} & seconds (s) \\
\bottomrule
\end{tabular}
\end{center}

\begin{tcolorbox}[warnbox]
The variable names are \textbf{case-sensitive}.
\texttt{R} must be in \textbf{meters} (not \textmu m) and \texttt{t}
must be in \textbf{seconds} (not \textmu s).
If your data is in other units, convert it in MATLAB before saving.
\end{tcolorbox}

\begin{figure}[H]
\centering
\includegraphics[width=\textwidth]{Fig/fig1_variable_name.png}
\caption{Example MATLAB workspace showing the required variable
  format: \texttt{R} (bubble radius, in meters) and \texttt{t}
  (time, in seconds). Both arrays must be one-dimensional and of
  equal length.}
\label{fig:format_variable}
\end{figure}

% ---- Step 2 -------------------------------------------------
\step{Launch the Software and Load Your Data}

Double-click \texttt{IMR\_GUI.exe} to open the program.

Once the window appears, go to \textbf{File} $\rightarrow$
\textbf{Load experiment data (.mat)}, and select your \texttt{.mat}
file.

\begin{figure}[H]
\centering
\includegraphics[width=\textwidth]{Fig/fig2.png}
\caption{Opening the experimental data file via
  \textbf{File} $\rightarrow$ \textbf{Load experiment data (.mat)}.}
\label{fig:load_data}
\end{figure}

After loading, the experimental bubble radius vs.\ time curve will
appear in the plot area on the right.

\begin{figure}[H]
\centering
\includegraphics[width=\textwidth]{Fig/fig3.png}
\caption{Main window after successfully loading the dataset.
  The experimental bubble radius--time trace is displayed in the
  plot area. Time is shown in \textmu s (horizontal axis) and
  normalized radius on the vertical axis.}
\label{fig:show_data}
\end{figure}

% ---- Step 3 -------------------------------------------------
\step{Select the Constitutive Model}

In the \textbf{Model} dropdown (top-left panel), select:
\[
  \textbf{NHKV\,(Rmax)}
\]

\begin{figure}[H]
\centering
\includegraphics[width=\textwidth]{Fig/fig4.png}
\caption{Selecting \textbf{NHKV\,(Rmax)} from the \textbf{Model}
  dropdown in the top-left control panel.}
\label{fig:constitutional_model}
\end{figure}

\begin{tcolorbox}[notebox]
NHKV\,(Rmax) is recommended for IMR because it starts the simulation
from the experimentally observed maximum radius $R_\mathrm{max}$,
which avoids the need to specify an initial bubble-wall velocity.
\end{tcolorbox}

% ---- Step 4 -------------------------------------------------
\step{Enter Initial Parameter Guesses and Preview the Simulation}

In the \textbf{Parameters} panel, enter the following initial values:

\begin{center}
\begin{tabular}{lll}
\toprule
\textbf{Parameter} & \textbf{Symbol} & \textbf{Initial value} \\
\midrule
Shear modulus & $G$   & $1\times10^{4}$~Pa \\
Viscosity     & $\mu$ & $0.1$~Pa$\cdot$s \\
\bottomrule
\end{tabular}
\end{center}

\begin{figure}[H]
\centering
\includegraphics[width=\textwidth]{Fig/fig5.png}
\caption{Initial parameter guesses entered in the \textbf{Parameters}
  panel: shear modulus $G = 1\times10^{4}$~Pa and
  viscosity $\mu = 0.1$~Pa$\cdot$s. The checkboxes on the left
  indicate which parameters will be optimized during fitting.}
\label{fig:initial_guess}
\end{figure}

Click \textbf{Simulate} to preview the simulated curve with the
current parameters.
You can try different values of $G$ and $\mu$ to explore how each
parameter affects the shape of the simulated curve before running the
automated fitting.

\begin{figure}[H]
\centering
\includegraphics[width=\textwidth]{Fig/fig6.png}
\caption{Simulation preview: the simulated curve (orange) is overlaid
  on the experimental data (blue) after clicking \textbf{Simulate}.
  The two curves do not need to match at this stage --- this is only
  a visual check of the initial guess.}
\label{fig:simulation}
\end{figure}

% ---- Step 4.5 -----------------------------------------------
\halfstep{Switch to the Fitting Module}

After previewing the simulation, switch from the \textbf{Simulate}
tab to the \textbf{Fit} tab in the left control panel. This tab
contains all the controls needed to configure and run the automated
fitting: the algorithm selector, parameter bounds, fit window, and
the \textbf{Fit} button.

\begin{figure}[H]
\centering
\includegraphics[width=\textwidth]{Fig/fig6.5.png}
\caption{Switching from the \textbf{Simulate} tab to the
  \textbf{Fit} tab in the left control panel. All fitting-related
  settings --- algorithm, parameter bounds, and fit window --- are
  located in this tab.}
\label{fig:fitting_module}
\end{figure}

% ---- Step 5 -------------------------------------------------
\step{Set the Fitting Algorithm and Parameter Bounds}

\textbf{Algorithm:} In the \textbf{Fit} tab, set the algorithm to
\textbf{Pattern Search}.

\begin{figure}[H]
\centering
\includegraphics[width=\textwidth]{Fig/fig7.png}
\caption{Selecting \textbf{Pattern Search} from the
  \textbf{Algorithm} dropdown in the \textbf{Fit} tab.}
\label{fig:pattern_search}
\end{figure}

\textbf{Bounds:} Set the parameter search bounds as follows:

\begin{center}
\begin{tabular}{lll}
\toprule
\textbf{Parameter} & \textbf{Lower bound (LB)} & \textbf{Upper bound (UB)} \\
\midrule
$G$    & $1\times10^{3}$~Pa    & $2\times10^{6}$~Pa \\
$\mu$  & $0.01$~Pa$\cdot$s    & $1$~Pa$\cdot$s \\
\bottomrule
\end{tabular}
\end{center}

\begin{figure}[H]
\centering
\includegraphics[width=\textwidth]{Fig/fig8.png}
\caption{Lower bound (LB) and upper bound (UB) entered for each
  parameter in the \textbf{Parameters} panel.
  The optimizer will only search within these ranges.}
\label{fig:set_bound}
\end{figure}

% ---- Step 6 -------------------------------------------------
\step{Adjust the Fit Window}

By default, fitting uses the entire time record. Restrict it to the
\textbf{first expansion--collapse cycle} only.

In the \textbf{Fit window} section, drag the boundary handles (or
type values directly) so that the shaded region covers only the first
cycle --- from the start of the data up to where the bubble reaches
its first minimum radius after collapse.

\begin{figure}[H]
\centering
\includegraphics[width=\textwidth]{Fig/fig9.png}
\caption{Fit window (shaded region) restricted to the first
  expansion--collapse cycle. The right boundary is placed just after
  the bubble reaches its first post-collapse minimum radius.}
\label{fig:fitting_window}
\end{figure}

\begin{tcolorbox}[notebox]
Fitting only the first cycle is sufficient and avoids the influence
of secondary oscillations, which are harder to model accurately.
\end{tcolorbox}

% ---- Step 7 -------------------------------------------------
\step{Run the Fitting}

Click the \textbf{Fit} button. The optimizer will automatically adjust
$G$ and $\mu$ to minimize the difference between the simulated and
experimental curves.

\begin{figure}[H]
\centering
\includegraphics[width=\textwidth]{Fig/fig10.png}
\caption{Clicking the \textbf{Fit} button to start the optimizer.
  The status indicator shows that fitting is in progress. No further
  user interaction is required until the optimizer finishes.}
\label{fig:start_fitting}
\end{figure}

\textbf{You do not need to do anything while fitting is in progress.}
The parameter values and the simulated curve will update automatically
when the optimizer finishes.

\begin{tcolorbox}[notebox]
Fitting typically takes \textbf{1--5 minutes} depending on your
computer. The window will remain responsive during this time.
\end{tcolorbox}

% ---- Step 8 -------------------------------------------------
\step{Evaluate the Results}

When fitting is complete, check the following two criteria:

\begin{enumerate}[leftmargin=1.5em]
  \item \textbf{Visual match:} The simulated curve (solid line) should
        closely follow the experimental data (markers) during the
        \emph{first expansion--collapse cycle}.
        A small discrepancy in later cycles is acceptable.

  \item \textbf{LSQ error:} In the \textbf{Output} panel, check the
        value labeled \texttt{LSQErr}. An acceptable fit has:
        \[
          \texttt{LSQErr} < 500
        \]
\end{enumerate}

\begin{figure}[H]
\centering
\includegraphics[width=\textwidth]{Fig/fig11.png}
\caption{Fitting complete: the simulated curve (orange) closely
  follows the experimental data (blue) through the first
  expansion--collapse cycle. The \textbf{Output} panel (bottom-left)
  reports the optimized parameters $G$ and $\mu$ and the final
  \texttt{LSQErr}.}
\label{fig:end_fitting}
\end{figure}

\begin{tcolorbox}[okbox]
If both criteria are satisfied --- good visual overlap on the first
cycle \textbf{and} \texttt{LSQErr} $< 500$ --- you have
obtained a successful fit. Record the reported $G$ and $\mu$ values.
\end{tcolorbox}

If the fit is poor (large error or obvious mismatch), try a different
initial guess and run again. If you are still unable to obtain a
satisfactory fit, consult your TA.

% ---- Step 9 -------------------------------------------------
\step{Save and Submit Your Results}

Go to \textbf{File} $\rightarrow$ \textbf{Export result (.mat)\ldots}
and save the output file.

\begin{figure}[H]
\centering
\includegraphics[width=\textwidth]{Fig/fig12.png}
\caption{Saving the results via
  \textbf{File} $\rightarrow$ \textbf{Export result (.mat)\ldots}.
  The exported file contains the fitted parameters, the simulated
  curve, and the original experimental data.}
\label{fig:export_result}
\end{figure}

The exported \texttt{.mat} file contains the fitted parameters, the
simulated curve, and the experimental data. \textbf{Submit this file}
as your result.

\vspace{10pt}
\hrule
\vspace{8pt}

% ---- Troubleshooting ----------------------------------------
\section*{Quick Troubleshooting}

\begin{center}
\begin{tabular}{p{5.5cm} p{8cm}}
\toprule
\textbf{Problem} & \textbf{What to do} \\
\midrule
Data does not load &
  Check that your \texttt{.mat} file has variables named exactly
  \texttt{R} (meters) and \texttt{t} (seconds). \\[4pt]
Simulated curve is completely wrong &
  Verify the model is set to \textbf{NHKV\,(Rmax)} and that the
  physics constants ($P_\infty$, $\rho$) match your experiment. \\[4pt]
\texttt{LSQErr} is very large after fitting &
  Try a different initial guess. If $G$ is near a bound, widen the
  bound and retry. \\[4pt]
Anything else &
  Ask your TA before changing any setting not covered in this guide. \\
\bottomrule
\end{tabular}
\end{center}

\vspace{8pt}
\begin{tcolorbox}[notebox, title={Default physics constants}]
Unless instructed otherwise, leave all physics constants at their
default values:
$P_\infty = 101{,}325$~Pa,\quad
$\rho = 998$~kg/m$^3$,\quad
$\gamma = 0.056$~N/m.
\end{tcolorbox}

\end{document}
